\chapter{Mathematics}

\section{Quadratic equations}
\begin{itemize}[nosep,leftmargin=*]
\item $ax^2+bx+c=0$, $a\ne0$, $D=b^2-4ac$:
\[
 x_{1,2}=\frac{-b\pm\sqrt D}{2a},\qquad
 x_1+x_2=-b/a,\quad x_1x_2=c/a.
\]
\item $D>0,=0,<0$: two, one, no real roots.
Vertex $(-b/(2a),-D/(4a))$: minimum if $a>0$, maximum if $a<0$.
\item $a=0$: linear. $a=b=0$: all $x$ if $c=0$, otherwise none.
\item Widen intermediates; verify integer bounds exactly.
\end{itemize}

\section{Integer division and alignment}
\kactlimport{IntDivision.h}
\kactlimport{ModAlign.h}

\section{Sums and progressions}
\begin{itemize}[nosep,leftmargin=*]
\item $n\ge0$; empty sum $=0$. Arithmetic: $a_i=a_0+id$.
\[
 \sum_{i=0}^{n-1}a_i=na_0+\frac{dn(n-1)}2.
\]
\item Geometric: $a_i=a_0r^i$.
\[
 \sum_{i=0}^{n-1}a_i=
 \begin{cases}a_0(1-r^n)/(1-r),&r\ne1,\\na_0,&r=1.\end{cases}
\]
Infinite sum $a_0/(1-r)$ requires $|r|<1$.
\item Power sums:
\[
 \begin{aligned}
 \sum_{i=1}^n i&=\frac{n(n+1)}2,\\
 \sum_{i=1}^n i^2&=\frac{n(n+1)(2n+1)}6,\\
 \sum_{i=1}^n i^3&=\left(\frac{n(n+1)}2\right)^2,\\
 \sum_{i=1}^n i^4&=\frac{n(n+1)(2n+1)(3n^2+3n-1)}{30}.
 \end{aligned}
\]
\item For integers $0\le a\le b$ and $c\ne1$:
\[
 c^a+c^{a+1}+\cdots+c^b=\frac{c^{b+1}-c^a}{c-1}.
\]
For $c=1$, the sum is $b-a+1$.
\item $S(n)=\sum_{i=0}^{n-1}f(i)$: sum on $[l,r)$ is $S(r)-S(l)$.
\item Cancel exact factors or widen before multiplying.
Division modulo $m$ requires an invertible denominator.
\end{itemize}

% Sums/Series formulas from the supplied original KACTL Mathematics.
\section{Series}
\begin{itemize}[nosep,leftmargin=*]
\item For all real x:
\[
 e^x=1+x+\frac{x^2}{2!}+\frac{x^3}{3!}+\cdots.
\]
\item $-1<x\le1$:
\[
 \ln(1+x)=x-\frac{x^2}{2}+\frac{x^3}{3}-\frac{x^4}{4}+\cdots.
\]
\item $-1\le x\le1$:
\[
 \sqrt{1+x}=1+\frac{x}{2}-\frac{x^2}{8}
 +\frac{x^3}{16}-\frac{5x^4}{128}+\cdots.
\]
\item For all real x (radians):
\[
 \begin{aligned}
 \sin x&=x-\frac{x^3}{3!}+\frac{x^5}{5!}-\frac{x^7}{7!}+\cdots,\\
 \cos x&=1-\frac{x^2}{2!}+\frac{x^4}{4!}-\frac{x^6}{6!}+\cdots.
 \end{aligned}
\]
\end{itemize}

\section{Exact arithmetic and overflow}
\begin{itemize}[nosep,leftmargin=*]
\item Signed $b$-bit range: $[-2^{b-1},2^{b-1}-1]$; $b=64,128$.
\item Promote before arithmetic: \texttt{i128(a)*b}.
Check intermediate bounds, including negating the signed minimum.
\end{itemize}
\kactlimport{Int128.h}
\begin{itemize}[nosep,leftmargin=*]
\item \texttt{cpp\_int}: arbitrary-precision integers and stream I/O;
cost depends on bit length. Range-check before fixed-width conversion.
\end{itemize}
\begin{lstlisting}[moredelim={**[s][\color{codeString}]{<}{>}}]
#include <boost/multiprecision/cpp_int.hpp>
using boost::multiprecision::cpp_int;
\end{lstlisting}

\section{Probability}
\begin{itemize}[nosep,leftmargin=*]
\item $P(A\mid B)=P(A\cap B)/P(B)$, requiring $P(B)>0$.
\item Independence: $P(A\cap B)=P(A)P(B)$.
Pairwise independence need not imply mutual independence.
\item Partition $(B_i)$, omitting $P(B_i)=0$:
\[
 \begin{aligned}
 P(A)&=\sum_i P(A\mid B_i)P(B_i),\\
 P(B_j\mid A)&=\frac{P(A\mid B_j)P(B_j)}{P(A)}
 \quad(P(A)>0).
 \end{aligned}
\]
\end{itemize}

\section{Variance and common distributions}
\begin{itemize}[nosep,leftmargin=*]
\item $\operatorname{Var}(X)=\mathbb E[X^2]-\mathbb E[X]^2$.
For independent variables, $\operatorname{Var}(\sum a_iX_i)=\sum a_i^2\operatorname{Var}(X_i)$.
\item Hypergeometric: sample n without replacement from N objects, K marked; $0\le n,K\le N$:
$P(X=k)=\binom Kk\binom{N-K}{n-k}/\binom Nn$; impossible choices give zero.
\item Poisson, rate $\lambda\ge0$: $P(X=k)=e^{-\lambda}\lambda^k/k!$,
$k\ge0$; mean and variance $\lambda$.
\item Uniform on $[a,b]$: mean $(a+b)/2$, variance $(b-a)^2/12$.
Exponential rate $\lambda>0$: $P(T>t)=e^{-\lambda t}$, $t\ge0$;
mean $1/\lambda$, variance $1/\lambda^2$.
\end{itemize}

\section{Expectation}
\begin{itemize}[nosep,leftmargin=*]
\item Discrete $X$: $\mathbb E[X]=\sum_x xP(X=x)$.
For finite expectations, independence is unnecessary:
\[
 \mathbb E\!\left[c+\sum_i a_iX_i\right]
 =c+\sum_i a_i\mathbb E[X_i].
\]
\item Indicator $I_A$: $\mathbb E[I_A]=P(A)$;
expected count $=\sum_i P(A_i)$.
\item Partition $(B_i)$:
$\mathbb E[X]=\sum_iP(B_i)\mathbb E[X\mid B_i]$;
omit zero-probability events.
\end{itemize}

\section{Binomial and geometric distributions}
\begin{itemize}[nosep,leftmargin=*]
\item Binomial: successes in $n\ge0$ independent trials,
fixed $0\le p\le1$.
\[
 P(X=k)=\binom nk p^k(1-p)^{n-k},\quad 0\le k\le n.
\]
$\mathbb E[X]=np$; $p=0,1$ are deterministic.
Sampling without replacement is generally not binomial.
\item Geometric: independent trials, fixed $0<p\le1$;
$T$ includes the first success.
\[
 P(T=k)=(1-p)^{k-1}p,\quad k\ge1;\qquad
 \mathbb E[T]=1/p.
\]
Failures $T-1$ have mean $(1-p)/p$.
$p=1$: $T=1$; $p=0$: infinite waiting time.
\end{itemize}

\kactlimport{Probability.h}

\section{Random processes on states}
\begin{itemize}[nosep,leftmargin=*]
\item State contains all relevant history.
$p_{st}=P(\text{next}=t\mid\text{current}=s)$;
$p_{st}\ge0$, $\sum_t p_{st}=1$.
\item Hitting target before failure:
$h=1$ at targets; $h=0$ at failures and states unable to reach a target;
otherwise $h_s=\sum_t p_{st}h_t$.
\item Expected steps: $E=0$ at terminals;
$E_s=1+\sum_t p_{st}E_t$.
For expected immediate cost $c_s$, replace 1 by $c_s$.
\item Self-loop, $p_{ss}<1$:
\[
 E_s=\frac{c_s+\sum_{t\ne s}p_{st}E_t}{1-p_{ss}}.
\]
For hitting probability, omit $c_s$.
\item DAG: reverse topological DP, $O(V+E)$.
Cycles: real Gaussian elimination, $O(V^3)$.
\item Finite chain: expected steps are finite iff every reachable state
can reach a terminal through positive-probability edges.
A reachable closed nonterminal class implies infinite expected time.
\end{itemize}
\section{Trigonometry and recurrences}
\begin{itemize}[nosep,leftmargin=*]
\item Angles in radians:
$\sin(u+v)=\sin u\cos v+\cos u\sin v$;
$\cos(u+v)=\cos u\cos v-\sin u\sin v$.
$a\cos x+b\sin x=\sqrt{a^2+b^2}\cos(x-\operatorname{atan2}(b,a))$.
\item For $s_n=\sum_{i=1}^k c_i s_{n-i}$ with $c_k\ne0$: polynomial
$x^k-\sum_i c_ix^{k-i}$. Distinct roots $r_i$ give $s_n=\sum_i d_ir_i^n$;
a root of multiplicity m contributes a polynomial of degree m-1 times $r^n$.
Prefer \texttt{LinearRecurrence.h} for modular computation.
\end{itemize}
\par
